\section{Discussion}

The proposed formulation and methodology provide a structured approach to cooperative caching under dynamic network conditions. Unlike approaches that assume a fixed helper topology, this framework evaluates cache placement under sampled perturbations. This matters in IoT edge systems because effective connectivity can change even when the physical AP infrastructure remains fixed: users move, association patterns shift, APs can become unavailable, and some cooperative links become less useful over time.

A key insight from the measured results is the trade-off between redundancy and accessibility. The global popularity baseline stores the same objects everywhere and therefore produces high redundancy with weak cooperative benefit. Local popularity avoids global duplication but ignores neighbors, so it misses opportunities for nearby APs to diversify their contents. The topology-aware policies perform better because they use the mobility graph to decide when duplication is useful and when diversity is more valuable.

The hybrid Greedy + DQN result should be interpreted carefully. The learning component does not replace the greedy algorithm; it improves it locally. This is a practical design choice. Pure reinforcement learning over the full placement matrix would require a large action space and many training episodes. By starting from a greedy placement, the DQN only needs to find beneficial replacements around an already reasonable solution. The accepted-action inference rule also prevents the refiner from making changes that increase the objective.

The relocation term is also important. Without it, a learning policy can improve access cost by frequently rewriting caches, which may be unrealistic for storage-limited or bandwidth-limited edge devices. In the final run, the hybrid policy changes only about 0.2 objects per active node relative to greedy at the nominal perturbation setting, yet it improves the objective from 2.633 to 2.414. This suggests that modest, targeted changes are enough to improve the warm start.

Several limitations remain. First, the included run uses a deterministic trace-compatible fallback because the official Dartmouth archive is available through DataPort rather than direct unauthenticated download. The implementation is ready for a manually downloaded trace in the same movement format, but the reported numbers should be viewed as reproducible simulation results rather than measurements on the raw Dartmouth archive. Second, content demand is synthetic because public WLAN movement traces do not include object request logs. Third, the DQN uses explicit graph features rather than a graph neural network, so it may not generalize as well across graph sizes as a learned graph encoder.

Future work should evaluate the pipeline on the raw Dartmouth movement files after DataPort access is available, add temporal train/validation/test splits, sweep Zipf skew and cache capacity, and compare against online LRU/LFU replacement baselines. A graph neural network or multi-agent policy could also replace the current vector DQN, especially for larger topologies where permutation-aware graph encoders would be more appropriate.

Overall, the framework provides a complete and reproducible foundation for studying adaptive edge caching. Its main value is not only the final objective improvement, but the connection between mobility-derived topology, explicit cost modeling, executable optimization, and reportable evaluation artifacts.
